% Source presentations for register row N-c-022.

\subsection{From declarative clauses to SLD search and finite failure}
\label{sec:p09-presentation-n-c-022}

\subsubsection{One relation, two modes, and three kinds of control}

Kowalski's procedural reading starts from clauses, not from a separately
specified input--output function.  For a concrete source-faithful example, use
\([\,]\) for the empty list and \([u\mid v]\) for
\(\operatorname{cons}(u,v)\), and let
\[
\begin{array}{rcl}
 \mathit{app}([\,],Y,Y)&\leftarrow&\\
 \mathit{app}([U\mid X],Y,[U\mid Z])
   &\leftarrow&\mathit{app}(X,Y,Z).
\end{array}
\]
These are universally quantified definite clauses defining one ternary
relation.  Calling the first two arguments \enquote{inputs} and the third an
\enquote{output} describes one use of that relation; it is not part of the
logical clause.

In the forward mode, the goal
\(\leftarrow\mathit{app}([a,b],[c],Z)\) has the refutation
\[
\begin{aligned}
 \leftarrow\mathit{app}([a,b],[c],Z)
 &\Longrightarrow
 \leftarrow\mathit{app}([b],[c],Z_1)
   &&[Z=[a\mid Z_1]],\\
 &\Longrightarrow
 \leftarrow\mathit{app}([\,],[c],Z_2)
   &&[Z_1=[b\mid Z_2]],\\
 &\Longrightarrow\Box
   &&[Z_2=[c]].
\end{aligned}
\]
The computed answer is \(Z=[a,b,c]\).  In the inverse mode, the same two
clauses enumerate every split of a fixed list.  The complete branching
calculation for
\(\leftarrow\mathit{app}(X,Y,[a,b])\) is
\[
\begin{array}{rcll}
 \leftarrow\mathit{app}(X,Y,[a,b])
 &\Longrightarrow&\Box
   &[X=[\,],\ Y=[a,b]],\\
 &\Longrightarrow&
   \leftarrow\mathit{app}(X_1,Y,[b])
   &[X=[a\mid X_1]],\\[2pt]
 \leftarrow\mathit{app}(X_1,Y,[b])
 &\Longrightarrow&\Box
   &[X_1=[\,],\ Y=[b]],\\
 &\Longrightarrow&
   \leftarrow\mathit{app}(X_2,Y,[\,])
   &[X_1=[b\mid X_2]],\\[2pt]
 \leftarrow\mathit{app}(X_2,Y,[\,])
 &\Longrightarrow&\Box
   &[X_2=[\,],\ Y=[\,]].
\end{array}
\]
At the last node the recursive head does not unify with an empty third list, so
it contributes no child to the complete tree.
Thus the three successful answers are
\[
 (X,Y)=([\,],[a,b]),\qquad([a],[b]),\qquad([a,b],[\,]).
\]
This is the relational use of variables as partially specified data emphasized
in the early procedural presentation.  The complete member computation in
Kowalski supplies the original branching precedent; the append relation is the
worked source example in the later semantic and SLD developments
\cite[printed pp.~570--573]{Kowalski1974}
\cite{VanEmdenKowalski1976,AptVanEmden1982}.

Three choices must remain separate.
\begin{enumerate}
\item An \emph{atom-selection rule} chooses one atom of the current goal.
\item A \emph{clause choice} chooses one unifying program-clause head and hence
  one child of the complete tree.
\item A \emph{tree-search strategy} decides which pending child to visit next.
\end{enumerate}
The first is the computation rule quantified over by the SLD completeness
theorem.  The second creates logical nondeterminism.  The third includes
breadth-first, fair interleaving, and Prolog's depth-first backtracking.
Kowalski's paper supplies the procedural reading and the computation-space
examples; it does not identify any one traversal with logical consequence.

\subsubsection{Herbrand interpretations and the immediate-consequence proof}

Fix a first-order signature with at least one ground term.  Its \emph{Herbrand
universe} is the set of all ground terms over that signature, and its
\emph{Herbrand base} is the set of all ground predicate atoms over that
universe.  A ground instance of a clause is obtained by replacing every
variable with a ground term.  Consequently a Herbrand interpretation is exactly
a subset \(I\) of the Herbrand base: a ground atom is true precisely when it
belongs to \(I\).

For a definite program \(P\), Definition~\ref{def:p09-definite-program} gives
\[
 T_P(I)=
 \{\,H\theta:
   (H\leftarrow B_1,\ldots,B_k)\in P,\ 
   \theta\text{ ground},\
   B_i\theta\in I\text{ for all }i\,\}.
\]
This is a semantic operator on sets of ground atoms.  It is not one SLD
transition: it applies every ground clause instance simultaneously.

The finite ancestor program makes the iteration visible:
\[
\begin{array}{ll}
 parent(a,b).&parent(b,c).\\
 anc(X,Y)\leftarrow parent(X,Y).&
 anc(X,Z)\leftarrow parent(X,Y),anc(Y,Z).
\end{array}
\]
Starting from the empty interpretation,
\[
\begin{aligned}
 T_P^1(\varnothing)
   &=\{parent(a,b),parent(b,c)\},\\
 T_P^2(\varnothing)
   &=T_P^1(\varnothing)
     \cup\{anc(a,b),anc(b,c)\},\\
 T_P^3(\varnothing)
   &=T_P^2(\varnothing)\cup\{anc(a,c)\},\\
 T_P^4(\varnothing)&=T_P^3(\varnothing).
\end{aligned}
\]
The consequence \(anc(a,c)\) appears independently of a clause order, selected
goal atom, or search traversal.

The least-model theorem follows from four explicit steps.

\paragraph{Monotonicity.}
If \(I\subseteq J\), every ground clause body contained in \(I\) is contained
in \(J\), so \(T_P(I)\subseteq T_P(J)\).

\paragraph{Models are exactly pre-fixed points.}
A Herbrand interpretation \(I\) satisfies every ground instance
\[
 H\leftarrow B_1,\ldots,B_k
\]
exactly when membership of all \(B_i\) in \(I\) implies membership of \(H\).
This is precisely \(T_P(I)\subseteq I\).

\paragraph{The union of the finite stages is a fixed point.}
Put
\[
 L=\bigcup_{n<\omega}T_P^n(\varnothing).
\]
Every clause body is finite.  If all its atoms belong to \(L\), there is one
finite stage containing all of them: take the maximum of their finitely many
stage indices.  Hence
\[
 T_P(L)
 =\bigcup_{n<\omega}T_P(T_P^n(\varnothing))
 =\bigcup_{n<\omega}T_P^{n+1}(\varnothing)
 =L.
\]

\paragraph{It is the least model.}
If \(I\) is any Herbrand model, then \(T_P(I)\subseteq I\).  Induction on \(n\)
gives \(T_P^n(\varnothing)\subseteq I\) for every \(n\), and therefore
\(L\subseteq I\).  Since \(L\) is itself a fixed point and hence a model, it is
the least Herbrand model:
\[
 M_P=L=\operatorname{lfp}(T_P).
\]
For a ground atom \(A\), ordinary first-order consequence agrees with
membership in this least model: by the Herbrand theorem, failure of ground
consequence has a Herbrand countermodel, while the intersection of the
Herbrand models of a definite program is itself its least model
\cite[printed pp.~736--741]{VanEmdenKowalski1976}.

The finite-stage statement has an equally concrete proof-tree form.  A
\emph{ground proof tree} for \(A\) has root \(A\); an internal node labelled
\(H\) chooses a ground clause instance
\(H\leftarrow B_1,\ldots,B_k\) and has the \(B_i\) as children; a fact is a
leaf.  Taking a fact-only tree to have height \(1\), induction on \(n\) proves
\[
 A\in T_P^n(\varnothing)
 \quad\Longleftrightarrow\quad
 A\text{ has a ground proof tree of height at most }n.
\]
The forward induction reads the clause instance that put \(A\) into the next
stage.  The reverse induction applies \(T_P\) after all child roots have
appeared.  Thus every member of \(M_P\) has a finite proof, even when the
program or Herbrand base is infinite; there need not be one uniform finite
height for all consequences.

\subsubsection{SLD answers: soundness, lifting, and strong completeness}

Now return to goal-directed proof.  Standardize every input clause apart before
unification.  For the goal \(\leftarrow anc(a,Z)\), choose the recursive clause
first:
\[
\begin{aligned}
 \leftarrow anc(a,Z)
 &\xRightarrow{\theta_1}
   \leftarrow parent(a,Y_0),anc(Y_0,Z),
 &\theta_1&=\{X_0/a,\ Z_0/Z\},\\
 &\xRightarrow{\theta_2}
   \leftarrow anc(b,Z),
 &\theta_2&=\{Y_0/b\},\\
 &\xRightarrow{\theta_3}
   \leftarrow parent(b,Z),
 &\theta_3&=\{X_1/b,\ Y_1/Z\},\\
 &\xRightarrow{\theta_4}\Box,
 &\theta_4&=\{Z/c\}.
\end{aligned}
\]
Here \((X_0,Y_0,Z_0)\) and \((X_1,Y_1)\) belong to fresh variants of the two
clauses.  Restricting the composed substitution to the initial variable gives
the computed answer \(Z=c\).  The root branch using the base clause instead
returns \(Z=b\).

\paragraph{Why every computed answer is correct.}
At one positive SLD step, the selected goal atom \(A_i\) is unified with the
head \(H\) of
\[
 H\leftarrow B_1,\ldots,B_k.
\]
The universally closed clause entails
\(B_1\land\cdots\land B_k\to H\).  After applying the mgu \(\theta\), the
conjunction of the new resolvent therefore entails the conjunction of the old
goal under \(\theta\).  Chaining these implications through a refutation and
using that the empty conjunction is true yields
\[
 P\models
 \forall\bigl((A_1\land\cdots\land A_n)
              \theta_1\cdots\theta_m\bigr).
\]
After restriction to the variables of the initial goal, this is exactly the
correct-answer condition in Definition~\ref{def:p09-sld}.  Apt and van Emden's
fixed-point proof sharpens the same invariant.  If \(G_{m-i}\) is the goal
\(i\) steps before success and \(\theta_{m-i+1},\ldots,\theta_m\) are the
remaining substitutions, then every ground instance of every atom in
\(G_{m-i}\theta_{m-i+1}\cdots\theta_m\) belongs to
\(T_P^i(\varnothing)\)
\cite[thm.~5.1 and cor.~5.2]{AptVanEmden1982}.

\paragraph{The lifting step.}
Suppose an instance \(G\eta\) takes an SLD step using a standardized-apart
clause \(C\) whose variables are disjoint from those of \(G\eta\), and write
\(\sigma\) for the mgu used in that step.  Then \(\eta\sigma\) unifies the
selected atom of \(G\) with the uninstantiated clause head.  If \(\theta\) is
an mgu of those two atoms, the defining property of an mgu gives a substitution
\(\rho\) through which this unifier factors:
\[
 (\eta\sigma)|_{\operatorname{Var}(G)\cup\operatorname{Var}(C)}
 =
 (\theta\rho)|_{\operatorname{Var}(G)\cup\operatorname{Var}(C)}.
\]
The resolvent from \(G\) under \(\theta\) therefore has the resolvent from
\(G\eta\) under \(\sigma\) as a further instance.  Repeating this argument
lifts every finite refutation of an instance to a refutation of the original
goal.  If \(\sigma_1,\ldots,\sigma_m\) are the mgus in the instance
refutation and \(\theta_1,\ldots,\theta_m\) those in the lifted refutation, the
general factorization is
\[
 \eta\sigma_1\cdots\sigma_m
   =(\theta_1\cdots\theta_m)\rho
\]
on the variables of the initial goal.  In the special case used below, the
instance refutation has the identity as its computed answer, and the left side
reduces there to \(\eta\).  This is why completeness says that every correct
answer is an \emph{instance} of a computed answer; the two substitutions need
not be literally equal.

\paragraph{From the least model to a refutation.}
For a ground atom \(A\in M_P\), choose the least stage at which \(A\) appears.
The corresponding ground clause has body atoms at earlier stages.  Induction
on that stage constructs finite refutations of those body goals.  The lifting
lemma replaces the ground clause instances by standardized-apart clauses and
mgus.  For a finite conjunction, the same induction combines the component
refutations.  Thus a semantically correct ground instance has an SLD
refutation.

For a non-ground correct answer \(\eta\), the universal closure of each
instantiated goal atom follows from \(P\).  Replace its remaining variables by
distinct constants occurring in neither \(P\) nor \(G\eta\).  Ground
completeness gives refutations for those ground atoms; replacing the fresh
constants back by variables gives
refutations of the atoms in \(G\eta\) whose computed answers are the identity.
Combine them into a refutation of \(G\eta\), then apply the lifting lemma.  The
result has computed answer \(\theta\) and a further substitution \(\rho\) such
that \(\eta=\theta\rho\) on the initial variables.  Finally, the switching lemma
permutes adjacent independent atom selections; repeated switching transforms a
refutation into one following any fixed computation rule, without changing the
answer except for variants.  This is the proof spine behind Lloyd's Theorems
9.2--9.5 and the exact strong answer theorem stated in
Theorem~\ref{thm:p09-sld-adequacy}
\cite[ch.~2, \S\S8--9]{Lloyd1987}.

\paragraph{What the theorem does not say.}
Consider the extensionally equivalent left-recursive presentation
\[
\begin{array}{rcl}
 anc(X,Y)&\leftarrow&anc(X,Z),parent(Z,Y),\\
 anc(X,Y)&\leftarrow&parent(X,Y),
\end{array}
\]
with the two parent facts above and the recursive clause placed first.  The
complete SLD tree for \(\leftarrow anc(a,Y)\) has a base-clause branch
\[
 \leftarrow parent(a,Y)\Longrightarrow\Box\quad[Y=b].
\]
Its first recursive branch begins
\[
\begin{aligned}
 \leftarrow anc(a,Y)
 &\Longrightarrow
 \leftarrow anc(a,Z_1),parent(Z_1,Y)\\
 &\Longrightarrow
 \leftarrow anc(a,Z_2),parent(Z_2,Z_1),parent(Z_1,Y)\\
 &\Longrightarrow\cdots .
\end{aligned}
\]
A leftmost, textual-clause-order, depth-first traversal remains on this infinite
branch and never visits the existing answer.  The computation rule selected the
leftmost atom; clause order chose the recursive child; depth-first traversal
caused starvation.  Strong completeness says that the complete tree for the
fixed atom-selection rule contains a refutation.  It does not say that an
unfair traversal finds it
\cite[\S6, discussion after thm.~6.6]{AptVanEmden1982}.

\subsubsection{Completion and a proof from a finitely failed tree}

Clark completion changes one-way program clauses into classical definitions.
For clauses \(p(\bar t_i)\leftarrow B_i\), with local variables renamed apart,
Definition~\ref{def:p09-clark-completion} gives
\[
 p(\bar x)\leftrightarrow
 \bigvee_i\exists\bar y_i\,
   (\bar x=\bar t_i\land B_i),
\]
together with the displayed free-term equality theory.  The equality axioms
matter: they make distinct constructors disjoint, make each constructor
injective, exclude a finite term from containing itself, and justify the mgu
cases in the finite-tree proof.  Metalinguistic unification and object-language
equality in the completion remain different things.

For the parent/ancestor program, the completed definitions simplify to
\[
\begin{aligned}
 parent(x,y)
 &\leftrightarrow
   ((x=a\land y=b)\lor(x=b\land y=c)),\\
 anc(x,z)
 &\leftrightarrow
   \bigl(parent(x,z)\lor
   \exists y\,(parent(x,y)\land anc(y,z))\bigr).
\end{aligned}
\]
The complete SLD tree for the ground goal
\(\leftarrow anc(c,a)\) is finite and has no success leaf:
\[
\begin{array}{c}
 \leftarrow anc(c,a)\\[2pt]
 \swarrow\hspace{45mm}\searrow\\[-2pt]
 \leftarrow parent(c,a)
 \hspace{20mm}
 \leftarrow parent(c,Y_1),anc(Y_1,a)\\
 \text{failure}\hspace{42mm}\text{failure at }parent(c,Y_1).
\end{array}
\]
The two root children come from the two ancestor clauses.  At each leaf, no
parent fact unifies with a first argument \(c\).

The completion proves the same negative conclusion.  Free-term equality gives
\(c\ne a\) and \(c\ne b\), so the parent biconditional entails
\(\forall y\,\neg parent(c,y)\).  Both disjuncts in the ancestor
biconditional are therefore false at \((c,a)\), whence
\[
 \operatorname{Comp}(P)\models\neg anc(c,a).
\]
This is not an appeal to elapsed search time: it uses every child of the root
and the fact that the whole tree is finite.

The general soundness proof has the same shape.  At a positive selected atom,
the completion and equality axioms entail the goal clause at a failed leaf when
there is no unifying head; at an internal node they entail the parent goal
clause from the conjunction of all child goal clauses.  Induction on tree depth
therefore propagates failure from every leaf to the root.  With negative
literals, an outer induction on the finite SLDNF rank handles subsidiary trees:
a subsidiary refutation proves its ground atom, while a subsidiary finite
failure proves its classical negation in the completion.  This is the
rank-then-depth proof of Lloyd's Theorem 15.4, whose positive core is Clark's
finite-failure argument
\cite{Clark1978}\cite[thm.~15.4]{Lloyd1987}.

Finiteness is load-bearing.  Add only the clause
\[
 loop(X)\leftarrow loop(X).
\]
The goal \(\leftarrow loop(a)\) has one infinite branch and neither a refutation
nor a finite failed tree.  Its completed definition is the tautology
\(\forall x\,(loop(x)\leftrightarrow loop(x))\), which entails neither
\(loop(a)\) nor \(\neg loop(a)\).
Thus \enquote{no answer has appeared} cannot be substituted for finite failure.

\subsubsection{Ground SLDNF, four outcomes, and the exact theorem boundary}

Use the normal program
\[
\begin{array}{rcl}
 bird(tweety)&\leftarrow&\\
 flies(X)&\leftarrow&bird(X),\mathsf{not}\ abnormal(X),\\
 loop(X)&\leftarrow&loop(X).
\end{array}
\]
There is no clause defining \(abnormal\).  The registered safe SLDNF rule
selects a negative literal only after its atom is ground.  Its mutually
recursive success and failure definitions are well-founded by the finite rank
from Definition~\ref{def:p09-sldnf}.

For \(\leftarrow flies(tweety)\), positive resolution first gives
\[
 \leftarrow bird(tweety),\mathsf{not}\ abnormal(tweety)
 \Longrightarrow
 \leftarrow\mathsf{not}\ abnormal(tweety).
\]
The subsidiary rank-\(0\) tree for
\(\leftarrow abnormal(tweety)\) is the one-node finite failed tree, so the
ground negative literal succeeds at rank \(1\), and the outer goal refutes.  No
variable binding is returned by the negative step; negation as failure is a
test.

The four operational outcomes can now be separated without euphemism:
\begin{center}
\begin{tabularx}{\textwidth}{@{}l X X@{}}
\toprule
Goal & Subsidiary behavior & Outer result \\
\midrule
\(\leftarrow\mathsf{not}\ abnormal(tweety)\)
  & The ground positive goal has a finite complete tree with no success.
  & Negative success; completion entails
    \(\neg abnormal(tweety)\). \\
\(\leftarrow\mathsf{not}\ bird(tweety)\)
  & The ground positive goal has a refutation using the fact.
  & Negative failure; no negative answer is returned. \\
\(\leftarrow\mathsf{not}\ loop(tweety)\)
  & The ground positive goal follows one infinite recursive branch.
  & Divergence; neither negative success nor negative failure. \\
\(\leftarrow\mathsf{not}\ parent(X,b)\)
  & The negative atom is nonground and no positive literal remains to bind it.
  & Floundering under safe selection; neither success nor failure. \\
\bottomrule
\end{tabularx}
\end{center}

Safeness and allowedness have different standings.  Safeness is a restriction
on which literal the computation rule may select.  Allowedness is the syntactic
condition in Definition~\ref{def:p09-normal-program}.  The flies clause above
is allowed because its variable occurs in the positive body atom \(bird(X)\).
The floundering query is not allowed because \(X\) occurs only in a negative
literal.  Lloyd's Proposition 15.1 says that if the whole program/query pair is
allowed, no computation flounders and every computed answer grounds the
original query variables.  This is a sufficient syntactic guarantee, not the
definition of SLDNF success and not a hypothesis of the two soundness theorems.

Theorem 15.4 says that a finitely failed safe SLDNF tree makes the universally
closed goal clause a consequence of \(\operatorname{Comp}(P)\).  Theorem 15.6
says that every safe SLDNF computed answer is a correct completion answer.  Its
proof inducts on refutation length; positive steps use the completed defining
biconditional and negative steps use Theorem 15.4.  Neither theorem claims
termination or completeness for every normal program
\cite[\S15, prop.~15.1 and thms.~15.4, 15.6]{Lloyd1987}.

The adjacent completeness theorem has much stronger hypotheses.  A
hierarchical program has a natural-number level for each predicate such that
every body predicate has strictly smaller level than its clause head.  This
excludes both \(anc\) recursion and \(loop\).  Under that hypothesis,
allowedness, and a safe computation rule, Lloyd's Theorem 16.3 proves that the
SLDNF tree exists and is finite and that every ground correct completion answer
is computed.  Its termination proof orders each goal by the finite multiset of
predicate levels; every child strictly decreases that well-founded multiset.
This theorem cannot be relabelled as completeness for arbitrary recursive
normal programs \cite[thm.~16.3]{Lloyd1987}.

The internal SLDNF locators were checked in a university-hosted scan of the
corrected printing against the Springer edition record; the scan is not a
publisher-hosted copy, so that provenance remains explicit.  Concrete Prolog
is commonly presented with leftmost selection, textual clause order, and
depth-first traversal; it also provides cut and arithmetic, and some
implementations use unification without the occurs check \cite{Lloyd1987}.
The confirmed ISO
general-core record and public WG17 material separately identify control
constructs, cut, arithmetic, negation as failure, ordinary unification, and
unification with an occurs check as subjects addressed separately.  The
corrected normative clauses are not part of the available source basis, so no
clause-exact ISO behavior is asserted.  In every case these implementation and
language-standard choices do not alter the pure SLD theorem; they fall outside it unless separately
formalized \cite{ISO13211-1-1995}.
